\chapter{Menopausal hormone therapy (MHT)}
\index{Menopausal hormone therapy (MHT)}

\section*{Definition and Overview}
Menopausal hormone therapy (MHT), formerly known as hormone replacement therapy, is a pharmacological treatment that involves the administration of synthetic or bioidentical oestrogens, often combined with progestogens, to alleviate symptoms associated with the decline of endogenous ovarian hormones during perimenopause and menopause. The primary indication for MHT is the relief of moderate-to-severe vasomotor symptoms (such as hot flushes and night sweats), genitourinary syndrome of menopause (including vaginal dryness, dyspareunia, and recurrent urinary tract infections), and the prevention of postmenopausal osteoporosis. MHT is tailored to the individual, using various systemic (oral tablets, transdermal patches, gels) or local (vaginal creams, pessaries) formulations. In modern medicine, the decision to initiate MHT involves a careful risk-benefit analysis based on the patient's age, time since menopause, medical history, and personal preferences.

\section*{Epidemiology}
Menopause is a universal biological milestone, with the average age of onset being 51 years. In many countries, approximately 80\% of women experience menopausal symptoms, with about 20--30\% describing these symptoms as severe and disruptive to their quality of life. The clinical use of MHT has fluctuated significantly over the past few decades. Following the publication of the Women's Health Initiative study in 2002, which raised concerns about risks such as breast cancer and cardiovascular events, MHT use declined dramatically. However, subsequent re-analyses and contemporary guidelines have refined these risk patterns, establishing that for healthy symptomatic women under the age of 60 years, or within 10 years of menopause onset (known as the "window of opportunity"), the benefits of MHT generally outweigh the risks.

\section*{Aetiology and Pathophysiology}
Menopause is defined as the permanent cessation of menstruation resulting from the loss of ovarian follicular activity. The decline in oestrogen production drives the pathophysiology of menopausal symptoms.
\begin{itemize}
    \item \textbf{Hypoestrogenism:} The key hormonal driver, causing thermoregulatory instability in the hypothalamus, atrophy of the vulvovaginal epithelium, and accelerated bone resorption.
    \item \textbf{Vasomotor instability:} Fluctuations and subsequent lack of oestrogen disrupt hypothalamic neurotransmitters (such as serotonin and noradrenaline), narrowing the thermoregulatory zone and triggering flushing and sweating.
    \item \textbf{Genitourinary atrophy:} Oestrogen receptors are highly concentrated in the lower urogenital tract. Loss of oestrogen reduces tissue elasticity, blood flow, and vaginal secretions, leading to vaginal pH changes and susceptibility to infections.
    \item \textbf{Bone resorption acceleration:} Oestrogen typically inhibits osteoclast activity; its deficiency leads to rapid bone mineral density loss, particularly in the first 5--10 years post-menopause.
    \item \textbf{Progestogen requirement:} Systemic oestrogen therapy in women with an intact uterus causes endometrial hyperplasia and increases the risk of endometrial carcinoma. Progestogens must be co-administered to protect the endometrium.
\end{itemize}

\section*{Clinical Features}
The features prompting MHT evaluation relate directly to the physical and psychological changes of menopause.
\begin{itemize}
    \item Vasomotor symptoms, presenting as sudden, intense sensations of heat in the face, neck, and chest, often accompanied by sweating, chills, and palpitations.
    \item Night sweats, which disrupt sleep, leading to chronic fatigue, insomnia, and secondary cognitive impairment (commonly described as brain fog).
    \item Genitourinary syndrome of menopause, characterised by vaginal dryness, burning, itching, pain during sexual intercourse, urinary frequency, and urgency.
    \item Neuropsychiatric changes, including mood swings, irritability, anxiety, and depressive symptoms.
    \item Physical signs such as joint and muscle pain, breast tenderness, skin dryness, and weight gain with abdominal fat redistribution.
    \item Unexplained bone fractures or a diagnosed decrease in bone mineral density on dual-energy X-ray absorptiometry scans.
\end{itemize}

\section*{Differential Diagnosis}
Before initiating MHT, other conditions causing similar symptoms must be ruled out:
\begin{itemize}
    \item \textbf{Thyroid disease:} Hyperthyroidism can cause heat intolerance, sweating, palpitations, and weight loss, mimicking vasomotor symptoms, but is distinguished by thyroid function blood tests.
    \item \textbf{Pheochromocytoma or carcinoid syndrome:} Rare endocrine tumours that cause flushing, sweating, and hypertension, which must be differentiated from menopausal flushes.
    \item \textbf{Primary sleep disorders:} Conditions like sleep apnoea that cause night sweats and insomnia, requiring careful sleep studies to distinguish from vasomotor-induced sleep disturbance.
    \item \textbf{Clinical depression:} Mood disorders that present during midlife, which may coexist with menopause but require distinct psychiatric management.
    \item \textbf{Premature ovarian insufficiency:} Menopause occurring before the age of 40 years; requires distinct, high-dose hormone replacement strategies until the natural age of menopause.
\end{itemize}

\section*{When to Seek Medical Care}
Women experiencing disruptive menopausal symptoms should consult their general practitioner or a specialist menopause clinic. Certain symptoms warrant immediate or urgent medical review, especially if the patient is taking MHT.

\textbf{Contact your local emergency services for:}
\begin{itemize}
    \item Sudden, severe chest pain, shortness of breath, or coughing up blood (suggestive of myocardial infarction or pulmonary embolism).
    \item Sudden weakness or numbness in the face, arm, or leg, difficulty speaking, or visual changes (suggestive of a stroke).
    \item Severe, unilateral swelling and pain in one calf or leg (suggestive of deep vein thrombosis).
\end{itemize}

Other non-emergency reasons to seek medical advice include:
\begin{itemize}
    \item Experiencing unexpected or persistent vaginal bleeding (breakthrough bleeding) while on MHT.
    \item Noticing a new lump, skin changes, or nipple retraction in either breast.
    \item Finding that your menopausal symptoms are not controlled by your current MHT dose.
    \item Experiencing bothersome side effects, such as severe breast tenderness, bloating, headaches, or mood changes.
    \item Wishing to discuss stopping MHT or reviewing your ongoing MHT safety.
\end{itemize}

\section*{Diagnosis and Investigations}
Initiating MHT requires a comprehensive medical evaluation and targeted screening.
\begin{itemize}
    \item \textbf{Clinical history and physical examination:} A detailed assessment of menopausal symptoms, cardiovascular risk factors, personal or family history of breast cancer, thromboembolism, or liver disease.
    \item \textbf{Blood pressure check:} Essential baseline and monitoring measurement, as uncontrolled hypertension is a contraindication to oral MHT.
    \item \textbf{Mammography and breast examination:} Up-to-date breast screening must be confirmed before starting MHT and continued annually or biennially during treatment.
    \item \textbf{Pelvic ultrasound:} Indicated if there is a history of abnormal vaginal bleeding to assess endometrial thickness and rule out endometrial pathology.
    \item \textbf{Hormonal testing:} Not routinely required in women over 45 years with typical symptoms, but essential in women under 40 to diagnose premature ovarian insufficiency (measuring FSH and oestradiol).
    \item \textbf{Bone mineral density scan:} Performed to assess baseline bone health if MHT is being initiated primarily for osteoporosis prevention in at-risk women.
\end{itemize}

\section*{Management}
MHT management must be highly individualised, utilising the lowest effective dose for the shortest duration necessary to achieve symptom control.
\begin{itemize}
    \item \textbf{Oestrogen-only MHT:} Prescribed solely for women who have undergone a hysterectomy, as there is no risk of endometrial hyperplasia.
    \item \textbf{Combined MHT:} Required for women with an intact uterus, incorporating a progestogen (such as micronised progesterone) alongside oestrogen to protect the endometrium.
    \item \textbf{Transdermal versus oral administration:} Transdermal formulations (patches, gels) bypass the liver, carrying a lower risk of deep vein thrombosis and stroke compared to oral tablets, making them preferred for women with cardiovascular risk factors.
    \item \textbf{Local vaginal oestrogen:} Low-dose creams, pessaries, or rings used exclusively for vaginal dryness and genitourinary symptoms; these have minimal systemic absorption and do not require progestogen protection.
    \item \textbf{Tibolone:} A synthetic steroid MHT alternative with oestrogenic, progestogenic, and androgenic properties, suitable for postmenopausal women.
    \item \textbf{Non-hormonal alternatives:} Selective serotonin reuptake inhibitors, serotonin-noradrenaline reuptake inhibitors, clonidine, or gabapentin for women who cannot take MHT (such as breast cancer survivors).
    \item \textbf{Regular clinical review:} Annual assessment of symptom control, treatment side effects, breast screening compliance, and ongoing indication for therapy.
\end{itemize}

\section*{Complications}
\begin{itemize}
    \item \textbf{Thromboembolic disease:} Oral oestrogen increases the risk of deep vein thrombosis and pulmonary embolism, especially in the first year of use.
    \item \textbf{Breast cancer risk:} Long-term combined MHT is associated with a small, duration-dependent increase in breast cancer risk; oestrogen-only MHT carries a lower or neutral risk.
    \item \textbf{Stroke and cardiovascular disease:} Oral MHT slightly increases stroke risk when initiated in women over 60 years or more than 10 years post-menopause.
    \item \textbf{Endometrial hyperplasia and cancer:} Occurs if oestrogen therapy is administered to a woman with a uterus without adequate progestogen co-therapy.
    \item \textbf{Minor side effects:} Bloating, breast tenderness, headaches, nausea, and breakthrough vaginal bleeding, which are often dose-dependent.
\end{itemize}

\section*{Prognosis and Outcomes}
MHT is highly effective, providing rapid and complete relief of vasomotor symptoms in over 85--90\% of women. It also successfully halts postmenopausal bone loss, significantly reducing the lifetime risk of osteoporotic fractures. Upon cessation of MHT, vasomotor symptoms can recur in approximately 50\% of women, regardless of whether the therapy is stopped abruptly or tapered gradually.

\section*{Prevention and Self-Care}
\begin{itemize}
    \item Apply transdermal patches or gels to clean, dry, intact skin on the lower abdomen, hips, or thighs; never apply them on or near the breasts.
    \item Maintain a healthy lifestyle with a balanced diet rich in calcium and vitamin D to support bone health alongside MHT.
    \item Engage in regular weight-bearing and resistance exercises to strengthen bones and improve cardiovascular fitness.
    \item Dress in layers, use fans, and keep your bedroom cool to help manage hot flushes without relying solely on medications.
    \item Limit consumption of trigger substances, such as caffeine, alcohol, spicy foods, and smoking, which can precipitate vasomotor flushes.
    \item Use over-the-counter vaginal lubricants (during intercourse) and moisturisers (regularly) in addition to, or instead of, local oestrogen for mild vaginal dryness.
    \item Attend all scheduled mammography screenings and perform regular breast self-checks while taking MHT.
\end{itemize}

\section*{Sources and Further Reading}
The following organisations provide authoritative, evidence-based information on menopausal hormone therapy and menopause management:
\begin{itemize}
    \item Australasian Menopause Society. \textit{MHT clinical guidelines, symptom management, and risk tables.} https://www.menopause.org.au/
    \item Jean Hailes for Women's Health. \textit{Menopause management, symptoms, and MHT resources.} https://www.jeanhailes.org.au/
    \item NHS. \textit{Hormone replacement therapy overview.} https://www.nhs.uk/conditions/hormone-replacement-therapy-hrt/
    \item Mayo Clinic. \textit{Hormone therapy: benefits and risks, what to expect.} https://www.mayoclinic.org/healthy-lifestyle/womens-health/in-depth/hormone-therapy/art-20046330
    \item International Menopause Society. \textit{Global consensus statements on MHT safety.} https://www.imsociety.org/
\end{itemize}
